\documentclass[11pt,a4paper]{article}
\usepackage[T1]{fontenc}
\usepackage[utf8]{inputenc}
\usepackage[british]{babel}
\usepackage{lmodern,microtype}
\usepackage[margin=29mm]{geometry}
\usepackage{amsmath,amssymb,amsthm,mathtools}
\usepackage{booktabs,graphicx}
\usepackage{xcolor}
\usepackage[colorlinks=true,linkcolor=blue!50!black,citecolor=blue!50!black,urlcolor=blue!50!black]{hyperref}
\numberwithin{equation}{section}
\newtheorem{theorem}{Theorem}[section]
\newtheorem{proposition}[theorem]{Proposition}
\newtheorem{lemma}[theorem]{Lemma}
\newtheorem{corollary}[theorem]{Corollary}
\theoremstyle{definition}
\newtheorem{definition}[theorem]{Definition}
\theoremstyle{remark}
\newtheorem{remark}[theorem]{Remark}
\newcommand{\R}{\mathbb R}
\newcommand{\N}{\mathbb N}
\newcommand{\E}{\mathbb E}
\newcommand{\PP}{\mathbb P}
\newcommand{\ind}{\mathbf 1}
\newcommand{\dd}{\,\mathrm d}
\newcommand{\vol}{\operatorname{vol}}
\newcommand{\diam}{\operatorname{diam}}
\newcommand{\Law}{\mathcal L}
\newcommand{\Mop}{\mathcal M}
\DeclareMathOperator{\adj}{adj}
\title{Exact reciprocal-simplex product densities\newline below a prescribed spatial scale}
\author{}
\date{}
\begin{document}
\maketitle

\begin{abstract}
Spatial earthquake statistics motivate point-process models whose second,
third and fourth product densities are proportional to the reciprocals of
pair distance, triangle area and tetrahedron volume. Exact identities of this
kind cannot hold at every scale for a stationary process of positive
intensity. We show that they can hold simultaneously on every nondegenerate
configuration below any prescribed diameter, at any independently prescribed
positive intensity. The construction chooses between two whole-space laws:
a process on separated random spheres and a process of separated clusters
containing at most three points. Sphere geometry supplies the exact fourth
product density, while the small clusters correct the two lower densities.
All factorial expectation measures through order four are absolutely
continuous on the whole configuration space, with densities finite almost
everywhere. The process is locally finite and has finite third count moments,
although its fourth count moment is infinite in every nonempty open ball.
The construction is nonergodic.
\end{abstract}

\section{Introduction}

Earthquake hypocentres form spatial patterns concentrated around faults and
fault systems. A point-process description asks how these patterns appear in
the frequencies of pairs, triples and larger collections of events. The
work of Kagan and Knopoff on the two-point correlation function
\cite{KaganKnopoff1980}, followed by Kagan's studies of the three-point and
four-point moment functions \cite{Kagan1981Three,Kagan1981Four}, connects
these frequencies with the geometry of the simplex determined by the
selected events. Their empirical comparisons suggest approximate reciprocal
distance, area and volume relations. In the three-point study, the
comparison is with a Poisson catalogue having matched boundaries and depth
distribution. The four-point study explicitly discusses a divergent
height integral and a lower-height regularisation near the base plane
\cite[p.~730]{Kagan1981Four}. These observations motivate a compatibility
question; they do not assert the exact identities that we study below.

Daley and Vere-Jones \cite[Exercise 12.6.11]{DVJ2008} discuss these forms in
a question about their compatibility with a stationary point process.
Their formulation uses proportional or asymptotic notation; it does not
supply the global equality assumptions considered here. We make the domain
and the meaning of the densities explicit. We ask for exact identities for
the \emph{full} ordered product densities, and we distinguish identities at
all distances from identities below a fixed spatial scale. This distinction
is necessary: a reciprocal-distance second product density at every
distance already contradicts the growth of the mean count in large balls.

Fix a positive intensity $\lambda$ and a positive diameter $M$. Our main
result constructs a stationary simple point process on $\R^3$ for which
\[
 \begin{aligned}
 \rho^{(2)}(x_1,x_2)&=\frac{a_2}{D(x_1,x_2)},\\
 \rho^{(3)}(x_1,x_2,x_3)&=\frac{a_3}{A(x_1,x_2,x_3)},\\
 \rho^{(4)}(x_1,x_2,x_3,x_4)&=\frac{a_4}{V(x_1,x_2,x_3,x_4)}
 \end{aligned}
\]
whenever the corresponding configuration is nondegenerate and has diameter
less than $M$, apart from Lebesgue-null sets. Here $D$, $A$ and $V$ are
ordinary distance, area and volume. The positive constants $a_2,a_3,a_4$
are supplied by the construction. They are not independently prescribed.
There is no lower distance, area, volume or altitude bound, so very thin
triangles and almost coplanar tetrahedra remain in the required domain.

The fourth formula has a consequence that is easy to miss. If a fourth
vertex approaches the plane of a fixed noncollinear triangle, its
reciprocal tetrahedron volume grows as the reciprocal of the normal
distance to that plane. Integration in that normal direction diverges.
Thus any process satisfying the local fourth identity must have infinite
fourth count moments on open balls. This does not prevent a count from
being finite almost surely. We use factorial expectation measures as
extended nonnegative measures, and prove separately that their densities
are finite almost everywhere.

The construction begins with a surprisingly simple description. Toss one
fair coin and use its outcome to choose the entire point-process law.
The first law places Poisson points on a stationary family of spheres
whose surfaces are separated by more than $M$. Four nearby points must
then lie on one sphere. Four affinely independent points determine a
unique circumsphere, and the Jacobian of that change of coordinates is
proportional to their tetrahedron volume. A suitable distribution of sphere
radii cancels the remaining factor, giving the exact fourth identity.

The same calculation gives the first three product densities of the
sphere process. The second and third have reciprocal-simplex factors
multiplied by Gaussian attenuation terms. An internal length parameter
makes their deficits small in the integrals needed for a correction,
without changing the intensity or the leading coefficients. The second
law uses clusters of at most three points to supply those deficits. Its
fourth product density vanishes on configurations of diameter below $M$.
A choice between whole laws then adds the desired product densities with
no mixed tuples between the two alternatives.

The spherical change of variables is an instance of the Euclidean
Blaschke--Petkantschin formula stated in Chapron
\cite[equation (1)]{Chapron2018}. We give the required specialisation and
the lower-order calculations. That geometric identity by itself does not
establish the existence of the process: the probability budget, the
correction kernels and the global regularity arguments are separate parts
of the proof. We make no novelty or priority claim for the geometric
identity or for the resulting construction.

The example is a realisability result for these exact local product
densities. It is nonergodic, and we do not claim that it gives a statistical
model fitted to earthquake catalogues. Its purpose is to separate the
compatibility question from additional requirements such as ergodicity,
finite fourth moments or an empirical fit.

\section{Factorial expectation measures and product densities}
\label{sec:preliminaries}

Let $X$ be a point process on $\R^3$, regarded as a random counting measure.
We write $X(B)$ for the number of points in a Borel set $B$. The process is
\emph{locally finite} if $X(B)<\infty$ almost surely for every bounded $B$,
and \emph{simple} if its points have multiplicity one. It is
\emph{stationary} if translating every point by any fixed vector leaves
its law unchanged. A stationary process has intensity $\lambda$ when
\begin{equation}
 \E X(B)=\lambda |B|
 \label{eq:intensity-def}
\end{equation}
for every bounded Borel set $B$, where $|B|$ is three-dimensional Lebesgue
measure. In this paper $0<\lambda<\infty$.

\subsection{Counting ordered tuples}

For $n\geq1$, the $n$th factorial expectation measure $\alpha_X^{(n)}$ is
defined by
\begin{equation}
 \int h\dd\alpha_X^{(n)}
 =\E\sum_{x_1,\ldots,x_n\in X}^{\ne} h(x_1,\ldots,x_n),
 \qquad h\geq0\text{ measurable}.
 \label{eq:factorial}
\end{equation}
The superscript $\ne$ means that the selected points are pairwise distinct,
and the tuples are ordered. For disjoint sets $B_1,\ldots,B_n$, this
measure gives $\E\prod_i X(B_i)$ on $B_1\times\cdots\times B_n$.
For a single set $B$, it gives
\begin{equation}
 \alpha_X^{(n)}(B^n)=\E\bigl(X(B)\bigr)_n,
 \qquad (m)_n=m(m-1)\cdots(m-n+1).
 \label{eq:falling}
\end{equation}
Both sides may be infinite. These are the factorial moment measures in
the terminology of \cite[Chapter 9]{DVJ2008}; ``expectation measure''
here emphasises that local finiteness of the measure is not assumed.

If there is a measurable $\rho_X^{(n)}\geq0$ such that
\begin{equation}
 \alpha_X^{(n)}(\dd x_1\cdots\dd x_n)
 =\rho_X^{(n)}(x_1,\ldots,x_n)\dd x_1\cdots\dd x_n,
 \label{eq:product-density}
\end{equation}
we call $\rho_X^{(n)}$ a \emph{product density} or \emph{factorial density}.
The equality is an identity of full extended measures. It includes all
ordered tuples and leaves no singular measure on collinear, coplanar or
other exceptional configurations. We require the density to be finite
Lebesgue-almost everywhere, but its integral over a bounded set may be
infinite. A finite-almost-everywhere function need not be locally
integrable.

At intensity $\lambda$, the normalised functions are
\begin{equation}
 g_n(x_1,\ldots,x_n)
 =\frac{\rho_X^{(n)}(x_1,\ldots,x_n)}{\lambda^n}.
 \label{eq:gn}
\end{equation}
A homogeneous Poisson process has $\rho^{(n)}=\lambda^n$ and $g_n=1$.
Normalisation changes only the coefficients in the reciprocal-simplex
identities. It does not turn a full product density into a cumulant or a
connected correlation function. We use full product densities throughout.

\subsection{Two elementary constructions}
\label{subsec:two-constructions}

We will use a Poisson process with a random intensity measure. If $\Lambda$
is a random locally finite nonatomic measure and, conditionally on
$\Lambda$, $Y$ is a Poisson process with intensity measure $\Lambda$, then
\begin{equation}
 \alpha_Y^{(n)}=\E\Lambda^{\otimes n}.
 \label{eq:cox-identity}
\end{equation}
To see this, first use the Poisson factorial identity for products of
pairwise disjoint sets, condition on $\Lambda$, and then extend to
nonnegative measurable functions by simple functions and monotone
convergence. The nonatomic assumption makes all product diagonals null.
Such a process is called a Cox process. Its directing measure may be
singular with respect to volume; in our example it is supported on sphere
surfaces. Identity \eqref{eq:cox-identity} is still valid.

A different operation is a mixture of laws. If a coin chooses $Y$ with
probability $\theta$ and $C$ with probability $1-\theta$, defining the whole
process $X$, then conditioning on that coin gives
\begin{equation}
 \alpha_X^{(n)}
 =\theta\alpha_Y^{(n)}+(1-\theta)\alpha_C^{(n)}.
 \label{eq:mixture-measure}
\end{equation}
This follows directly from \eqref{eq:factorial}, also when expectations
are infinite. A realisation contains points from the selected law only.
This differs from superposing two independent processes, whose full
product densities also contain tuples using points from both processes.

\section{The diameter-capped problem}
\label{sec:problem}

For points in $\R^3$, put
\begin{align}
 D(x_1,x_2)&=|x_2-x_1|,\nonumber\\
 A(x_1,x_2,x_3)&=\frac12|(x_2-x_1)\times(x_3-x_1)|,\label{eq:simplex}\\
 V(x_1,x_2,x_3,x_4)&=\frac16
 \left|\det[x_2-x_1,x_3-x_1,x_4-x_1]\right|.\nonumber
\end{align}
Write $(\Delta_2,\Delta_3,\Delta_4)=(D,A,V)$ and define
\begin{equation}
 E_n(M)=\{(x_1,\ldots,x_n):\diam\{x_1,\ldots,x_n\}<M,
                       \ \Delta_n(x_1,\ldots,x_n)>0\}.
 \label{eq:domains}
\end{equation}
The strict diameter bound does not bound the simplex away from degeneracy.
For example, triangles in $E_3(M)$ can have arbitrarily large circumradius.

\begin{theorem}\label{thm:main}
For every independently prescribed pair $\lambda>0$ and $M>0$, there is a
stationary, simple, locally finite point process $X$ on $\R^3$ with intensity
$\lambda$, and finite positive constants $a_2,a_3,a_4$, such that the full
factorial expectation measures of orders $2,3,4$ have globally defined
Lebesgue densities finite almost everywhere and
\begin{equation}
 \rho_X^{(n)}(\mathbf x)=\frac{a_n}{\Delta_n(\mathbf x)}
 \quad\text{for almost every }\mathbf x\in E_n(M),\qquad n=2,3,4.
 \label{eq:target}
\end{equation}
The constants and the process may depend on $(\lambda,M)$.
The process has finite third count moments on every bounded set and
infinite fourth count moments in every nonempty open ball. It is
nonergodic under translations.
\end{theorem}

There is no claim that a prescribed triple $(a_2,a_3,a_4)$ can be realised.
Outside the sets $E_n(M)$ the densities are unconstrained, subject to the
global measure and finiteness requirements of the theorem. Values of a
density on a Lebesgue-null set may be chosen arbitrarily; this permission
does not permit a singular component of its measure.

For clarity, let us distinguish this problem, denoted S1, from the
unrestricted version S0, in which the three equalities are required on all
nondegenerate configurations without a diameter cap. S0 is false.

\begin{proposition}\label{prop:global-obstruction}
No stationary point process of finite positive intensity has full second
factorial density $a_2/|x_2-x_1|$ almost everywhere on distinct pairs in
$\R^3$, with $0<a_2<\infty$ and no other part of the second factorial
measure.
\end{proposition}
\begin{proof}
Let $B_R$ be a ball of radius $R$ and set $N_R=X(B_R)$. The assumed density
is integrable on $B_R^2$, since
\[
 \int_{B_R}\int_{B_R}\frac{\dd y\dd x}{|x-y|}
 \leq |B_R|\int_{B(0,2R)}\frac{\dd u}{|u|}
 =8\pi R^2|B_R|.
\]
Thus $N_R$ has finite second moment, and nonnegativity of its variance gives
\[
 \lambda^2|B_R|^2
 \leq \E N_R^2
 =\lambda|B_R|+\E(N_R)_2
 \leq \lambda|B_R|+8\pi a_2R^2|B_R|.
\]
After dividing by $|B_R|^2=\bigl(4\pi R^3/3\bigr)^2$, we obtain
$\lambda^2\leq3\lambda/(4\pi R^3)+6a_2/R$, which is impossible as
$R\to\infty$.
\end{proof}

This obstruction uses no fourth-moment assumption. It requires a change
to the unrestricted tail of the second density, but it does not uniquely
determine a replacement problem. The common diameter cap in S1 is the
particular domain considered in Theorem~\ref{thm:main}.

\section{Existence by a mixture of two processes}
\label{sec:construction}

We begin with the final random choice, since it explains the job of each
component. Let $Y$ be a process on separated spheres, and let $C$ be a
process of separated clusters containing at most three points. The process
we seek has law
\begin{equation}
 \Law(X)=\frac12\Law(Y)+\frac12\Law(C).
 \label{eq:fair-coin}
\end{equation}
The coin is tossed once for the whole of $\R^3$. Our task is therefore to
construct two actual laws whose intensities average to $\lambda$ and whose
full product densities average to \eqref{eq:target}.

The correction process has a bounded intensity capacity because its
clusters are separated. We allow the sphere process to supply the rest.
Set
\begin{equation}
 \kappa=\frac1{36e\pi M^3},\qquad
 \lambda_C=\min\{\lambda,\kappa/2\},\qquad
 \lambda_Y=2\lambda-\lambda_C.
 \label{eq:component-intensities}
\end{equation}
Both intensities are strictly positive and their average is $\lambda$.
The construction below works for arbitrarily large $\lambda_Y$, while the
correction uses $\lambda_C\leq\kappa/2$.

\subsection{Process 1: Poisson points on separated spheres}
\label{subsec:spheres}

Let
\begin{equation}
 C_0=\frac{\pi^{3/2}}{3\sqrt3}.
 \label{eq:C0}
\end{equation}
Choose a finite $K>0$ sufficiently large that, with
\begin{equation}
 c=\sqrt{\frac{C_0K^3}{\lambda_Y}},\qquad
 F=\frac{K^4}{c^3}
  =\frac{\lambda_Y^{3/2}}{C_0^{3/2}\sqrt K},
 \qquad FM\leq10^{-4}.
 \label{eq:Kc}
\end{equation}
For example, $K=10^8M^2\lambda_Y^3/C_0^3$ satisfies the last inequality
with equality. The parameter $L\geq M$ is left free for now. Define
\begin{equation}
 f(q)=\frac KL e^{-q/L^2},\qquad
 a(R)=\frac{c}{f(R^2)},\qquad
 \gamma(R)=\frac{2R f(R^2)^4}{c^3}
          =\frac{2FR}{L^4}e^{-4R^2/L^2}.
 \label{eq:sphere-weights}
\end{equation}
We will construct a stationary random family of spheres whose first
intensity on centre--radius space is $\gamma(R)\dd z\dd R$. On a sphere
with centre $z$ and radius $R$, the Poisson surface intensity will be
$a(R)/(2R)$. These choices make the fourth-order weight constant after
putting $q=R^2$; the calculation appears in
Section~\ref{subsec:sphere-densities}.

The required spheres cannot be chosen independently, because nearby
points would then also come from different spheres. Instead, choose a
single radius band and arrange its surfaces far enough apart. Put
\begin{equation}
 H=2M,\qquad b_k=2^kM,\qquad P_k=4b_k+2M,\qquad k=0,1,\ldots,
 \label{eq:band-parameters}
\end{equation}
and partition the positive radii into
\[
 I_0=(0,M],\qquad I_k=(b_k/2,b_k]\quad(k\geq1).
\]
Define the band probabilities by
\begin{equation}
 p_k=HP_k^3\sup_{R\in I_k}\gamma(R).
 \label{eq:band-probabilities}
\end{equation}
They obey $\sum_kp_k<0.9$ for every $L\geq M$, as we verify in
Lemma~\ref{lem:budget}. With probability $p_k$, select band $k$; with the
remaining probability, select an empty sphere family.

Conditionally on selecting band $k$, perform the following steps.
Choose $Z$ uniformly on $[0,P_k)^3$ and use the centres
$Z+P_km$, $m\in\mathbb Z^3$. Independently choose $U$ uniformly on
$[0,H)$ and, at each centre, use the candidate radii
$R=U+jH\in I_k$, $j\in\mathbb Z$. Retain each candidate sphere,
independently conditionally on $(k,Z,U)$, with probability
\begin{equation}
 \theta_k(R)=\frac{\gamma(R)HP_k^3}{p_k}\in[0,1].
 \label{eq:sphere-retention}
\end{equation}
Each realised band contains only finitely many candidate radial indices
at each centre. A radius zero, which has phase probability zero, is
omitted. Independent marks on the countable set of centre and radius
indices specify this random family.

\begin{lemma}[Probability budget]\label{lem:budget}
The probabilities in \eqref{eq:band-probabilities} satisfy
\[
 \sum_{k\geq0}p_k<9000FM\leq0.9
\]
uniformly over $L\geq M$.
\end{lemma}
\begin{proof}
Since $P_k\leq6b_k$, the first band satisfies
\[
 p_0\leq864FM(M/L)^4\leq864FM.
\]
For $k\geq1$, use $R\leq b_k$ and $R>b_k/2$ in
\eqref{eq:sphere-weights}. With $t_k=b_k/L$, this gives
\[
 p_k\leq864FM\,t_k^4e^{-t_k^2}.
\]
The $t_k$ form a dyadic sequence. For the terms $t_k\leq1$, the sum of
$t_k^4$ is at most $16/15$. For $t>1$, the bound
$e^{t^2}\geq t^6/6$ gives $t^4e^{-t^2}\leq6/t^2$; the sum over the
remaining dyadic terms is at most $8$. Therefore
\[
 \sum_kp_k\leq864FM\left(1+\frac{16}{15}+8\right)
 =\frac{130464}{15}FM<9000FM.
\]
The bound contains no further dependence on $L$.
\end{proof}

\begin{lemma}[Sphere family]\label{lem:sphere-family}
The retained sphere family is stationary, its surfaces are pairwise
separated by more than $M$, and, for every nonnegative measurable $h$,
\begin{equation}
 \E\sum_{(z,R)\text{ retained}}h(z,R)
 =\int_{\R^3}\int_0^\infty h(z,R)\gamma(R)\dd R\dd z.
 \label{eq:sphere-first-intensity}
\end{equation}
\end{lemma}
\begin{proof}
Uniform translation of the lattice gives, for nonnegative $g$,
\[
 \E_Z\sum_{m\in\mathbb Z^3}g(Z+P_km)=P_k^{-3}\int_{\R^3}g(z)\dd z.
\]
Uniform radial phase similarly gives
\[
 \E_U\sum_{j:\,U+jH\in I_k}g(U+jH)
 =H^{-1}\int_{I_k}g(R)\dd R.
\]
Multiplying these identities by $p_k\theta_k(R)$, and summing over the
bands, proves \eqref{eq:sphere-first-intensity}. Only the first intensity
is asserted; the sphere parameters are dependent.

At a common centre the surfaces of two candidate spheres have distance
at least $H=2M$. At different centres in band $k$, their distance is at
least
\begin{equation}
 P_k-2b_k=2b_k+2M>M.
 \label{eq:surface-gap}
\end{equation}
These bounds survive retention. Finally, translating space changes $Z$
modulo $P_k\mathbb Z^3$ and relabels the centre indices. The conditional
retention law is unchanged by that relabelling, while the law of $Z$ is
uniform. This proves stationarity within each band and for their mixture.
\end{proof}

Write $S(z,R)$ for the sphere and $\sigma_{z,R}$ for its unnormalised
surface area measure. Given the retained family, place independent Poisson
processes on its spheres, with intensity measures
\begin{equation}
 \frac{a(R)}{2R}\sigma_{z,R}.
 \label{eq:surface-intensity}
\end{equation}
Equivalently, sample a Poisson number with mean $2\pi a(R)R$ on each
retained sphere and sample its directions independently and uniformly on
$S^2$. Their union is the process $Y$. Its directing measure is
\begin{equation}
 \Lambda=\sum_{(z,R)\text{ retained}}
                   \frac{a(R)}{2R}\sigma_{z,R}.
 \label{eq:Lambda}
\end{equation}

The construction already gives an actual stationary, simple and locally
finite process. To check local finiteness, condition on the selected band.
Only centres within distance $b_k$ of a bounded observation set can affect
it, and there are finitely many such centres and candidate radii. Each
sphere carries a finite Poisson count. The surface measures have no atoms,
and different surfaces are disjoint, so the union is simple. Countable
exhaustion by balls makes these statements simultaneous over the whole
space. The spatially invariant Poisson draws preserve stationarity.

We will compute the local product densities first. The fact that these
are globally well-defined densities requires attention to tuples lying
on several spheres, even though such tuples do not occur in $E_n(M)$.
The full proof is in Appendix~\ref{app:regularity}. It establishes the
following facts, which will also be used at the end of the construction.

\begin{proposition}\label{prop:sphere-regularity}
For the process $Y$, the full factorial expectation measures through
order four are absolutely continuous on the whole configuration space
and have densities finite almost everywhere. Moreover,
$\E Y(B)^3<\infty$ for every bounded Borel set $B$.
\end{proposition}

\subsection{Sphere geometry and the first four densities}
\label{subsec:sphere-densities}

Four noncoplanar points determine a unique sphere. The corresponding
change of coordinates is particularly useful here. Let $\omega_i\in S^2$
and $x_i=z+R\omega_i$, $i=1,\ldots,4$. The Euclidean spherical formula
stated in \cite[equation (1)]{Chapron2018}, specialised to dimension three,
is
\begin{equation}
 \dd x_1\cdots\dd x_4
 =6\Delta_3(\omega_1,\ldots,\omega_4)R^8
       \dd z\dd R\prod_{i=1}^4\dd\omega_i,
 \label{eq:chapron}
\end{equation}
where $\dd\omega$ is surface area measure on the unit sphere and
$\Delta_3$ is ordinary tetrahedron volume. Since
$V(x_1,\ldots,x_4)=R^3\Delta_3(\omega_1,\ldots,\omega_4)$, the Jacobian
is $6VR^5$. Chapron attributes the Euclidean identity to Miles; we use
the displayed formula in Chapron, without making a separate attribution
claim based on the original Miles paper.

A family with first intensity $\gamma(R)\dd z\dd R$, carrying Poisson
surface intensity $t(R)\sigma_{z,R}$, therefore has same-sphere fourth
density
\begin{equation}
 \frac{\gamma(R_*)t(R_*)^4R_*^3}{6V},
 \label{eq:generic-fourth}
\end{equation}
evaluated at the unique circumradius $R_*$. This follows by dividing
the parameter weight $\gamma t^4R^8$ by the Jacobian $6VR^5$.
The formula is ordered: no division by $4!$ is present.

For our process, separation makes this same-sphere contribution the full
local fourth density. We also need the lower orders, for which the set of
possible sphere centres has positive dimension. The next lemma records
all four calculations in a common convention.

\begin{lemma}[Sphere integrals]\label{lem:sphere-integrals}
Let $w_n:(0,\infty)\to[0,\infty]$ be measurable. Consider the measure
obtained by integrating
\[
 w_n(q)\prod_{i=1}^n\frac{\sigma_{z,\sqrt q}(\dd x_i)}{2\sqrt q}
\]
over $z\in\R^3$ and $q>0$. On nondegenerate configurations its density is
\begin{align}
 r_1&=2\pi\int_0^\infty\sqrt q\,w_1(q)\dd q,\label{eq:sphere-lemma1}\\
 r_2(x_1,x_2)&=\frac{\pi}{2D}
                 \int_{D^2/4}^\infty w_2(q)\dd q,\label{eq:sphere-lemma2}\\
 r_3(x_1,x_2,x_3)&=\frac1{8A}
                 \int_{\R}w_3(R_\triangle^2+h^2)\dd h,\label{eq:sphere-lemma3}\\
 r_4(x_1,x_2,x_3,x_4)&=\frac{w_4(R_*^2)}{48V}.\label{eq:sphere-lemma4}
\end{align}
Here $R_\triangle$ is the triangle circumradius and $R_*$ is the
tetrahedron circumradius. The formulas hold as nonnegative measure
identities; local integrability is not assumed.
\end{lemma}
\begin{proof}
The one-dimensional radial coarea formula identifies
\[
 \frac{\sigma_{z,\sqrt q}(\dd x)}{2\sqrt q}
 =\delta(|x-z|^2-q)\dd x.
\]
We use the right-hand side only as notation for the measure on the left.
Thus the density is obtained by integrating
\begin{equation}
 w_n(q)\prod_{i=1}^n\delta(|x_i-z|^2-q)
 \label{eq:incidence}
\end{equation}
over $(z,q)$. We justify this representation after computing the
constraint Jacobians.

For one point, translating $x_1$ to zero and using polar coordinates in
$z$ gives
$\int_{\R^3}\delta(|z|^2-q)\dd z=2\pi\sqrt q$.
For two points, write $x_1=0$ and $x_2=de$, with $|e|=1$. Subtracting the
first constraint from the second gives $d^2-2d(z\cdot e)=0$.
It fixes $z\cdot e=d/2$ with Jacobian $2d$. The first constraint fixes
$q=d^2/4+|z_\perp|^2$ with Jacobian one. Integrating the remaining two
centre coordinates in polar form gives
\[
 \frac1{2d}\int_{\R^2}w_2(d^2/4+|v|^2)\dd v
 =\frac\pi{2d}\int_{d^2/4}^\infty w_2(q)\dd q.
\]

For three points, translate the first to zero and write the others as
$u,v$. The two subtracted constraints fix the projection of the centre
onto the triangle plane at its circumcentre. Their Jacobian in orthonormal
plane coordinates is $4|u\times v|=8A$. If $h$ is the remaining normal
coordinate, the first constraint fixes $q=R_\triangle^2+h^2$.
This gives \eqref{eq:sphere-lemma3}.

For four noncoplanar points, the four constraints determine $(z,q)$
uniquely, with $q=R_*^2>0$. Subtracting the first row from the other three
in the derivative with respect to $(z,q)$ gives absolute determinant
\begin{equation}
 8\left|\det[x_2-x_1,x_3-x_1,x_4-x_1]\right|=48V.
 \label{eq:48V}
\end{equation}
This proves \eqref{eq:sphere-lemma4}, consistently with
\eqref{eq:chapron} after the change $q=R^2$.

For completeness, these calculations can be read as ordinary changes of
variables on the incidence set $|x_i-z|^2=q$. Work first in a compact
regular coordinate chart, with $q$ bounded away from zero and with a
continuous test function. Replace each displayed $\delta$ by a smooth
one-dimensional approximate identity. Integrating first in the radial
coordinate of each $x_i$ gives the surface measure divided by $2\sqrt q$.
Alternatively, keep the $x_i$ fixed and use $q$ and $n-1$ centre coordinates
as the constraint coordinates. The nonzero determinants just computed
give \eqref{eq:sphere-lemma1}--\eqref{eq:sphere-lemma4} in the limit. The
approximate-identity passage is the usual one-dimensional substitution
on these regular charts, with compact support and continuous Jacobians.
A countable chart decomposition proves equality of the resulting Borel
measures, and nonnegative exhaustion extends it to arbitrary measurable
weights and tests. All subsequent integrations use Tonelli's theorem.
No product of unspecified distributions, or finite fourth moment, is
needed.
\end{proof}

For $Y$, every tuple of diameter below $M$ lies on a single retained
sphere. By \eqref{eq:cox-identity} and
\eqref{eq:sphere-first-intensity}, its full local factorial measure is
therefore
\begin{equation}
 \int_{\R^3}\int_0^\infty
 \gamma(R)\left(\frac{a(R)}{2R}\right)^n
       \sigma_{z,R}^{\otimes n}(\dd\mathbf x)\dd R\dd z.
 \label{eq:local-sphere-measure}
\end{equation}
For $n=1$ the same expression gives the full first measure everywhere.
Putting $q=R^2$, the weight in Lemma~\ref{lem:sphere-integrals} becomes
\begin{equation}
 w_n(q)=\frac{\gamma(\sqrt q)a(\sqrt q)^n}{2\sqrt q},
 \qquad (w_1,w_2,w_3,w_4)=\left(\frac{f^3}{c^2},\frac{f^2}{c},f,c\right).
 \label{eq:wn}
\end{equation}
Define
\begin{equation}
 A_2=\frac{\pi K^2}{4c},\qquad A_3=\frac{K\sqrt\pi}{8}.
 \label{eq:A2A3}
\end{equation}
The four densities are
\begin{align}
 \rho_Y^{(1)}&=\frac{2\pi}{c^2}\int_0^\infty\sqrt q\,f(q)^3\dd q
              =\frac{C_0K^3}{c^2}=\lambda_Y,
              \label{eq:Y1}\\
 \rho_Y^{(2)}&=\frac{A_2}{D}e^{-D^2/(2L^2)}
                         &&\text{on }E_2(M),\label{eq:Y2}\\
 \rho_Y^{(3)}&=\frac{A_3}{A}e^{-R_\triangle^2/L^2}
                         &&\text{on }E_3(M),\label{eq:Y3}\\
 \rho_Y^{(4)}&=\frac{c}{48V}
                         &&\text{on }E_4(M).\label{eq:Y4}
\end{align}
For the first line,
\[
 \int_0^\infty q^{1/2}e^{-3q/L^2}\dd q
 =L^3\sqrt\pi/(6\sqrt3).
\]
The pair integral uses
\[
 \int_{D^2/4}^\infty e^{-2q/L^2}\dd q
 =(L^2/2)e^{-D^2/(2L^2)},
\]
and the triangle integral uses
$\int_\R e^{-h^2/L^2}\dd h=L\sqrt\pi$.
Thus $\lambda_Y,A_2,A_3,c$ are independent of the still-free scale $L$.

The fourth density is already exactly of the desired form. After the
fair coin it will have coefficient $c/96$. The lower two densities fall
short of the reciprocal forms because of their exponential factors.
Their pointwise deficits are nonnegative. We need more than pointwise
smallness: the triangle deficit contributes pairs when realised by
three-point clusters, so its pair marginal must be controlled uniformly,
including for thin triangles. That is the next step.

\subsection{The two deficits and their marginal bound}
\label{subsec:deficits}

For a translation-invariant kernel $h$ of $j$ points, write
\[
 (\Mop h)(x_1,\ldots,x_{j-1})
 =\int_{\R^3}h(x_1,\ldots,x_j)\dd x_j.
\]
Its \emph{anchored integral} is the integral with the first point fixed at
zero and the other points integrated over $\R^3$. Fixing one point removes
the infinite factor arising from translation invariance.
Set $K_2=\ind_{E_2(M)}/D$, with value zero on degenerate pairs, and put
\begin{equation}
 B_L=\frac{8\pi A_3M^3}{L},\qquad \delta=2B_L.
 \label{eq:BL}
\end{equation}
Define the desired correction densities
\begin{align}
 q_3&=\frac{A_3\ind_{E_3(M)}}{A}
                \left(1-e^{-R_\triangle^2/L^2}\right),\label{eq:q3}\\
 q_2&=\frac{\ind_{E_2(M)}}{D}
                \left[\delta+A_2\left(1-e^{-D^2/(2L^2)}\right)\right].
                \label{eq:q2}
\end{align}
Both kernels are symmetric, translation invariant and nonnegative, and
we set them to zero on degenerate configurations. The extra term
$\delta/D$ in the pair correction allows for pairs that must accompany a
three-point cluster. It changes the eventual coefficient $a_2$, which the
problem leaves free.

\begin{lemma}[Uniform marginal bound]\label{lem:deficit}
For every $L\geq M$,
\begin{equation}
 \Mop q_3\leq B_LK_2.
 \label{eq:marginal-bound}
\end{equation}
Consequently, the kernels
\begin{equation}
 F_3=q_3,\qquad F_2=q_2-\Mop q_3
 \label{eq:F-kernels}
\end{equation}
are nonnegative, with $F_2\geq B_LK_2$. Their anchored integrals are
finite and strictly positive.
\end{lemma}
\begin{proof}
Fix a pair $0,de$, with $0<d<M$ and $|e|=1$. For the third vertex $z$,
let $r$ be its distance from the line through the pair. The triangle area
is $dr/2$, and its circumradius is
\begin{equation}
 R_\triangle=\frac{|z|\,|z-de|}{2r}\leq\frac{M^2}{2r}
 \quad\text{on }E_3(M).
 \label{eq:circumradius-bound}
\end{equation}
Let $\varepsilon=M^2/(2L)\leq M/2$. Since
$1-e^{-t}\leq\min\{1,t\}$ for $t\geq0$, we have
\begin{align*}
 d\int_{\R^3}q_3(0,de,z)\dd z
 &\leq2A_3\int_{B(0,M)}\frac1r
                    \min\{1,\varepsilon^2/r^2\}\dd z\\
 &\leq8\pi A_3M\int_0^M\min\{1,\varepsilon^2/r^2\}\dd r\\
 &\leq16\pi A_3M\varepsilon=B_L.
\end{align*}
In the second line, cylindrical coordinates cancel the factor $1/r$,
and the axial length is at most $2M$. For the last line, split the
integral at $\varepsilon$; it is at most $2\varepsilon$. Collinear third
vertices form a null set. If the pair lies outside $E_2(M)$, no triangle
in the support of $q_3$ contains it, so the bound also holds there.

Let $Q_j$ be the anchored integral of $q_j$ and $Z_j$ that of $F_j$.
The pair integral $\int K_2(0,u)\dd u=2\pi M^2$, together with
$1-e^{-d^2/(2L^2)}\leq d^2/(2L^2)$, gives
\begin{equation}
 Q_2\leq\frac{32\pi^2A_3M^5}{L}
              +\frac{\pi A_2M^4}{2L^2}.
 \label{eq:Q2-bound}
\end{equation}
Indeed, the extra term integrates to $2\pi M^2\delta$, and the other
term is bounded by
\[
 4\pi A_2\int_0^M d^3/(2L^2)\dd d.
\]
Integrating \eqref{eq:marginal-bound} proves $Q_3<\infty$.
No infinite quantities are subtracted in \eqref{eq:F-kernels}.
The lower bound $q_2\geq2B_LK_2$ yields
$F_2\geq B_LK_2$, and Tonelli gives
\begin{equation}
 Z_3=Q_3>0,\qquad Z_2=Q_2-Q_3>0.
 \label{eq:ZQ}
\end{equation}
The strict positivity of $Q_3$ follows from the strictly positive
circumradius on every nondegenerate triangle; that of $Z_2$ follows from
its displayed lower bound.
\end{proof}

We can now choose the sphere scale. Take
\begin{equation}
 L=\max\left\{
 M,\ \frac{128\pi^2A_3M^5}{\lambda_C},\
 M^2\sqrt{\frac{2\pi A_2}{\lambda_C}}
 \right\}.
 \label{eq:L-choice}
\end{equation}
Each of the two summands in \eqref{eq:Q2-bound} is then at most
$\lambda_C/4$, so
\begin{equation}
 Q_2\leq\lambda_C/2.
 \label{eq:small-budget}
\end{equation}
All quantities on the right of \eqref{eq:L-choice} were fixed before
$L$. Lemma~\ref{lem:budget} remains valid at this value because its
probability bound was uniform in $L\geq M$.

The triangle deficit need not be uniformly small over shapes:
$R_\triangle$ is unbounded on $E_3(M)$. The estimate we use is instead the
uniform pair-marginal bound \eqref{eq:marginal-bound}. This is why no lower
altitude cutoff enters the argument.

\subsection{Process 2: separated finite correction clusters}
\label{subsec:correction}

We next realise $q_2$ and $q_3$ as the full local product densities of one
process with no eligible four-point tuples. Let $\Pi$ be a homogeneous
Poisson process of centres with intensity
\begin{equation}
 \eta=\frac1{36\pi M^3}.
 \label{eq:eta}
\end{equation}
Retain a centre if there is no other centre within distance $3M$. The
retained centres are dependent, but have separation greater than $3M$
almost surely and intensity
\begin{equation}
 \eta\exp\{-\eta|B(0,3M)|\}=\eta e^{-1}=\kappa.
 \label{eq:hardcore-intensity}
\end{equation}
Here $|B(0,3M)|=36\pi M^3$.

The first-intensity calculation can be obtained directly by Poisson
insertion. For a nonnegative measurable test $h$,
\begin{equation}
 \E\sum_{z\in\Pi}h(z,\Pi\setminus\{z\})
 =\eta\int_{\R^3}\E h(z,\Pi)\dd z.
 \label{eq:insertion}
\end{equation}
On a bounded region, condition on the Poisson count, select one point,
and use $n/n!=1/(n-1)!$ in the resulting nonnegative series.
Exhaustion proves \eqref{eq:insertion} in the whole space. Applying it
to the empty-ball condition gives \eqref{eq:hardcore-intensity} by the
Poisson void probability. Pairs at exactly distance $3M$ have zero
expected count in bounded windows, so strict separation holds almost
surely.

Set
\begin{equation}
 \beta=\lambda_C-Q_2+\frac{Q_3}{2}
       =\lambda_C-Z_2-\frac{Z_3}{2}.
 \label{eq:beta}
\end{equation}
By \eqref{eq:small-budget}, $\beta\geq\lambda_C/2>0$.
Independently at each retained centre, choose a cluster size in
$\{0,1,2,3\}$ with probabilities
\begin{equation}
 \pi_1=\frac\beta\kappa,\qquad
 \pi_2=\frac{Z_2}{2\kappa},\qquad
 \pi_3=\frac{Z_3}{6\kappa},\qquad
 \pi_0=1-\pi_1-\pi_2-\pi_3.
 \label{eq:cluster-probabilities}
\end{equation}
These are valid probabilities because
\begin{equation}
 \kappa(\pi_1+\pi_2+\pi_3)
 =\lambda_C-\frac{Q_2}{2}+\frac{Q_3}{6}
 \leq\lambda_C\leq\kappa.
 \label{eq:cluster-budget}
\end{equation}
We used $0<Q_3<Q_2$ in the inequality.

A size-zero cluster emits nothing and a size-one cluster emits its centre.
For size $j=2,3$, emit the centre $z$ and $j-1$ offsets
$z+y_2,\ldots,z+y_j$, with joint density
\begin{equation}
 \frac{F_j(0,y_2,\ldots,y_j)}{Z_j}
 \quad\text{on }(\R^3)^{j-1}.
 \label{eq:offset-law}
\end{equation}
Lemma~\ref{lem:deficit} proves that these are probability densities.
Every cluster has diameter below $M$ and contains its centre. The union
of the emitted points is $C$.

Points from different clusters are at distance greater than
$3M-2M=M$. Thus every tuple in $E_n(M)$ comes from one cluster. The
process is stationary because the thinning rule commutes with
translations and the offset laws do not depend on the centre location.
It is simple and locally finite: only Poisson centres within distance
$M$ of a bounded observation set can contribute, each emits at most
three points, and the continuous offset laws give no coincidences.

\begin{proposition}[The four correction densities]\label{prop:C-densities}
The process $C$ has intensity $\lambda_C$ and, on the indicated local
domains, full product densities
\begin{align}
 \rho_C^{(1)}&=\lambda_C,\label{eq:C1}\\
 \rho_C^{(2)}&=F_2+\Mop F_3=q_2
                      &&\text{on }E_2(M),\label{eq:C2}\\
 \rho_C^{(3)}&=F_3=q_3
                      &&\text{on }E_3(M),\label{eq:C3}\\
 \rho_C^{(4)}&=0
                      &&\text{on }E_4(M).\label{eq:C4}
\end{align}
\begingroup\emergencystretch=1em
Its full factorial expectation measures through order four have globally
finite-almost-everywhere densities, and its bounded-set count moments are
finite at every order.\par
\endgroup
\end{proposition}
\begin{proof}
Consider a cluster of size $j$ and an ordered selection of $n\leq j$ of
its points. After integrating its centre over space, symmetry and
translation invariance of $F_j$ make every ordered selection have the
same density. More explicitly, permute the selected point labels, choose
one selected point as the new anchor, and subtract its offset from the
remaining offsets. This linear change of variables has absolute
Jacobian one. The kernel is unchanged by the permutation and common
translation. Integrating the $j-n$ unused coordinates gives
$\Mop^{j-n}F_j$. There are $j!/(j-n)!$ ordered selections.

The retained-centre intensity $\kappa$ is sufficient for this calculation.
For pairs, size-two clusters contribute
$2\kappa\pi_2F_2/Z_2=F_2$, and size-three clusters contribute
$6\kappa\pi_3\Mop F_3/Z_3=\Mop F_3$. For triples, the contribution is
$6\kappa\pi_3F_3/Z_3=F_3$. There is no ordered selection of four points
from one cluster. Since separation excludes other local allocations,
these are the full local densities in \eqref{eq:C2}--\eqref{eq:C4}.
The first density is
\[
 \kappa(\pi_1+2\pi_2+3\pi_3)
 =\beta+Z_2+Z_3/2=\lambda_C.
\]
This calculation does not treat the retained centres as independent.

We still need the full measures outside the local domains. Give the same
independent marks to every original Poisson centre before retention, and
let $C^\sharp$ be the unthinned emitted process. Then $C\subseteq C^\sharp$
and, for bounded $B$,
\begin{equation}
 C^\sharp(B)\leq3\Pi(B^{+M}),
 \qquad B^{+M}=\{x:\operatorname{dist}(x,B)<M\}.
 \label{eq:C-domination}
\end{equation}
The right-hand side has finite moments of all orders. Absolute continuity
can be checked by retaining every parent allocation in the unthinned
process. Put $F_1=Z_1=1$, and for $1\leq n\leq3$ define the same-parent
kernel
\begin{equation}
 G_n=\eta\sum_{j=n}^3
        \frac{\pi_jj!}{(j-n)!Z_j}\Mop^{j-n}F_j,
 \qquad G_n=0\quad(n>3).
 \label{eq:G-kernels}
\end{equation}
In particular $G_1=\eta\sum_{j=1}^3j\pi_j$ is a constant.
Label a tuple by the partition $\mathcal P$ of its indices according to
its distinct Poisson parents. The Poisson factorial identity for those
parents, and the ordered-selection calculation above, give
\begin{equation}
 \rho_{C^\sharp}^{(n)}(\mathbf x)
 =\sum_{\mathcal P}\prod_{I\in\mathcal P}G_{|I|}(x_I).
 \label{eq:all-parent}
\end{equation}
The sum is over every set partition of $\{1,\ldots,n\}$. The identity
follows from iterated insertion \eqref{eq:insertion}; all summands are
nonnegative, so Tonelli applies. This proves global absolute continuity,
including all different-parent tuples. The moment bound
\eqref{eq:C-domination} makes the full measure locally finite, so its
density is finite almost everywhere.

For nonnegative tuple tests, the coupling gives
$\alpha_C^{(n)}\leq\alpha_{C^\sharp}^{(n)}$. Hence $C$ inherits global
absolute continuity and finite-almost-everywhere densities. It also
inherits the moment bounds from \eqref{eq:C-domination}.
\end{proof}

\subsection{The final mixture}

All the ingredients now exist at finite parameter values. The final step
is short because the two components were designed for the same coin.

\begin{corollary}\label{cor:mixture}
The law \eqref{eq:fair-coin} satisfies Theorem~\ref{thm:main}, with
\begin{equation}
 a_2=\frac{A_2+\delta}{2},\qquad
 a_3=\frac{A_3}{2},\qquad
 a_4=\frac{c}{96}.
 \label{eq:final-constants}
\end{equation}
Its normalised functions on $E_n(M)$ are
$g_n=a_n/(\lambda^n\Delta_n)$.
\end{corollary}
\begin{proof}
Both alternatives are stationary, simple and locally finite, and their
mean intensities average to $\lambda$ by
\eqref{eq:component-intensities}. Proposition~\ref{prop:sphere-regularity}
and Proposition~\ref{prop:C-densities} give globally finite-almost-everywhere
full densities. A finite mixture preserves these properties, and
\eqref{eq:mixture-measure}, \eqref{eq:Y2}--\eqref{eq:Y4} and
\eqref{eq:C2}--\eqref{eq:C4} give exactly \eqref{eq:final-constants}.

The parameter choices have no circular dependence. First
$(\lambda,M)$ determines $(\kappa,\lambda_C,\lambda_Y)$. Next choose
$K,c$ by \eqref{eq:Kc}, then obtain $A_2,A_3$ from
\eqref{eq:A2A3}. Choose $L$ from \eqref{eq:L-choice}, and finally define
$B_L,\delta$, the sphere band law and the cluster kernels. The uniform
sphere budget and \eqref{eq:small-budget} verify both probability laws
at that $L$. This works for every prescribed pair, without a restriction
on $\lambda M^3$.

Finite third count moments follow from the corresponding component
bounds. The fourth-moment assertion is proved just below. Finally, the
sphere law is empty with probability at least $1-\sum_kp_k>0.1$, so
$\PP(X=\varnothing)>0$. Positive intensity implies
$\PP(X\ne\varnothing)>0$. The event that the whole configuration is empty
is measurable and invariant under every translation, with probability
strictly between zero and one. Thus $X$ is nonergodic.
\end{proof}

\begin{proposition}\label{prop:fourth-infinite}
Any process with the full local fourth density $a_4/V$ on $E_4(M)$,
where $a_4>0$, has infinite fourth count moment in every nonempty open
ball.
\end{proposition}
\begin{proof}
Inside an arbitrary nonempty open ball, choose a smaller ball $B$ whose
diameter is less than $M$. Choose three small neighbourhoods of the
vertices of a fixed noncollinear triangle in its interior. Shrink them
so that every resulting base triangle has area bounded above and away
from zero and remains a positive distance from the boundary of $B$.
They form a positive-measure family of base triangles.

For each such triangle, a small planar disc about its barycentre and a
uniformly short normal interval remain inside $B$. Write the fourth
vertex in these coordinates as $(v,h)$. The tetrahedron volume is
$V=A|h|/3$, while the coordinate Jacobian is one. Apart from $h=0$, all
these configurations are in $E_4(M)$. The integral of $a_4/V$ over the
fourth vertex contains a positive multiple of
$\int_0^\varepsilon\dd h/h=\infty$. Tonelli, followed by integration
over the positive-measure family of base triangles, gives
$\alpha_X^{(4)}(B^4)=\infty$. Since $(X(B))_4\leq X(B)^4$, the raw
fourth moment is infinite as well.
\end{proof}

The same argument excludes the additional requirement of finite fourth
moments even in the capped problem. It does not exclude local finiteness,
and it does not impose any further geometric cutoff on the process we
have constructed.

% BEGIN SIMULATION SECTION
\section{Simulation study}
\label{sec:simulation}
\emph{Draft insertion point.} The reproducible simulation study is being
prepared separately. This section will state the sampling law, numerical
parameter choices, estimands, bin definitions, uncertainty estimates and
actual results for $g_2,g_3,g_4$ once those outputs have been checked. No
simulation result or empirical agreement is asserted in this draft.
% END SIMULATION SECTION

\appendix

\section{Global densities and moments of the sphere process}
\label{app:regularity}

This appendix proves Proposition~\ref{prop:sphere-regularity}. The local
formulas do not suffice for that proposition: they describe tuples on
one sphere in the prescribed diameter range, whereas the full measures
also include tuples spread across several spheres and several centres.
There are two distinct questions. We first prove absolute continuity for
every sphere allocation. We then show that the resulting sum of densities
is finite almost everywhere. At order four this second step needs a
geometric argument, because the fourth measure is not locally finite.

\subsection{Every allocation has an absolutely continuous measure}
\label{app:submersion}

Conditionally on the sphere family, $Y$ is Poisson with nonatomic locally
finite directing measure \eqref{eq:Lambda}. Hence
$\alpha_Y^{(n)}=\E\Lambda^{\otimes n}$ as extended measures. Fix a band
$k$ and a prescribed allocation of $n$ point labels to sphere indices
$(m_i,j_i)$, where $m_i\in\mathbb Z^3$ indexes a centre and
$j_i\in\mathbb Z$ indexes a radius. Repeated sphere indices are allowed.
A point in this allocation has the form
\begin{equation}
 x_i=Z+P_km_i+(U+j_iH)\omega_i,
 \qquad\omega_i\in S^2,\qquad i=1,\ldots,n.
 \label{eq:allocation-map}
\end{equation}
The parameter measure is a nonnegative measurable weight times Lebesgue
measure in $(Z,U)$ and product surface area measure in the directions.
Band membership, Poisson surface factors and the conditional retention
probabilities are part of this weight. When a sphere index is repeated,
its retention indicator is counted only once; the weight uses the
retention probability for each distinct selected sphere. This changes
no rank calculation below.

At positive radii, the angular derivatives of
\eqref{eq:allocation-map} span the two-dimensional plane
$\omega_i^\perp$ in the $i$th position block. Together these give $2n$
independent tangent directions. Quotienting them out, the derivative with
respect to $(Z,U)$ has $n$ rows
\begin{equation}
 (\omega_i^{\mathsf T},1),\qquad i=1,\ldots,n.
 \label{eq:normal-rows}
\end{equation}
For $n\leq4$ these rows have rank $n$ for almost every collection of
directions under product area measure. At order two, rank failure
requires the directions to coincide. Three distinct points on $S^2$
cannot be collinear, since a line intersects the sphere in at most two
points. Thus three directions are almost surely affinely independent.
Their affine plane intersects $S^2$ in a circle or a smaller set, of
surface measure zero, so a fourth independently integrated direction
almost surely avoids it. This proves the rank assertion, including all
lower orders.

It follows that the map \eqref{eq:allocation-map} has differential rank
$3n$ almost everywhere. Here is the resulting absolute-continuity argument
in detail. Use ordinary angular charts on the sphere. At a full-rank
parameter point, select a nonzero $3n$-by-$3n$ minor of the derivative.
The map consisting of the $3n$ position coordinates and the unused
parameter coordinates is locally invertible. In these coordinates the
position map is a coordinate projection. Fubini's theorem shows that the
preimage of any Lebesgue-null position set has parameter measure zero.
A countable cover by such inverse-function charts handles the full-rank
set; the rank-failure set was already null for the parameter measure.
Multiplication by a nonnegative measurable weight does not alter this
conclusion. Thus the pushforward for the specified allocation is
absolutely continuous.

There are countably many sphere allocations. Summing their measures, and
then mixing over the countably many bands, proves global absolute
continuity of $\alpha_Y^{(n)}$ for $n\leq4$. In particular, no mass remains
on a coincidence, collinearity or coplanarity set.

We can also exhibit the density without invoking a Radon--Nikodym theorem
for a potentially non-locally-finite fourth measure. For each fixed band
and bounded observation set, the number of influencing spheres has a
deterministic finite bound. Moreover,
\[
 2\pi a(R)R\leq2\pi b_k\sup_{0<R\leq b_k}a(R)<\infty.
\]
The directing mass on that set is therefore bounded conditionally on the
band, uniformly over the shift, phase and retention marks. All its count
moments are finite. The conditional factorial measures for a fixed band
are locally finite and absolutely continuous, so have finite-almost-everywhere
densities $h_{n,k}$. The full measure has the extended density
\begin{equation}
 \rho_Y^{(n)}=\sum_{k\geq0}p_kh_{n,k}.
 \label{eq:band-density-sum}
\end{equation}
The empty alternative contributes zero. A countable sum of finite densities
can still be infinite on a set of positive measure, so the remaining
parts of this appendix prove that this does not occur.

\subsection{Finite third count moments}
\label{app:third-moment}

Fix a ball $B=B(x,r)$. There is a constant $J_B<\infty$, independent of the
selected band, bounding the number of candidate spheres that can meet
$B$. Indeed, their centres must lie in $B(x,b_k+r)$, and the lattice
spacing satisfies $P_k\geq4b_k$ with $b_k\geq M$. A cubical packing bound
therefore gives a uniform finite bound on the number of such centres,
depending only on $r/M$. At any one centre, radii of spheres meeting $B$
lie in an interval of length at most $2r$. A grid of spacing $H=2M$
contains at most $\lfloor2r/H\rfloor+2$ such radii. Multiplying the two
bounds gives $J_B$; retention can only decrease the number.

We also have the uniform area bound
\begin{equation}
 \sigma_{z,R}(B)\leq C_B:=4\pi r^2.
 \label{eq:cap-bound}
\end{equation}
For $R\leq r$, use the full sphere area. For $R>r$, a nonempty
intersection is a spherical cap. If $d=|z-x|>0$, the cosine at its boundary
is
\[
 \frac{R^2+d^2-r^2}{2Rd}
 =\frac{d+(R^2-r^2)/d}{2R}
 \geq\frac{\sqrt{R^2-r^2}}{R}.
\]
Thus the cap area is at most
\[
 2\pi R^2\left(1-\sqrt{1-r^2/R^2}\right)\leq2\pi r^2.
\]
If $d=0$ and $R>r$, the intersection is empty, so the bound holds there
as well.

Write $m_{z,R}=a(R)\sigma_{z,R}(B)/(2R)$ for a sphere's contribution to
$\Lambda(B)$. The number bound and \eqref{eq:cap-bound} imply
\[
 \Lambda(B)^3\leq J_B^2\sum_{(z,R)\text{ retained}}m_{z,R}^3,
 \qquad
 \sigma_{z,R}(B)^3\leq C_B^2\sigma_{z,R}(B).
\]
Translation of the centre gives
\begin{equation}
 \int_{\R^3}\sigma_{z,R}(B)\dd z=4\pi R^2|B|.
 \label{eq:surface-centre-integral}
\end{equation}
Apply \eqref{eq:sphere-first-intensity},
\eqref{eq:surface-centre-integral} and Tonelli to obtain
\begin{align}
 \E\Lambda(B)^3
 &\leq\frac{\pi J_B^2C_B^2|B|}{2}
             \int_0^\infty\frac{\gamma(R)a(R)^3}{R}\dd R\nonumber\\
 &=\frac{\pi J_B^2C_B^2|B|}{2}
             \,2\int_0^\infty f(R^2)\dd R
 =\frac{\pi J_B^2C_B^2|B|}{2}K\sqrt\pi<\infty.
 \label{eq:third-Lambda}
\end{align}
The cancellation $\gamma a^3/R=2f(R^2)$ treats both small and large
radii. For a Poisson variable of mean $s$, the third raw moment is
$s^3+3s^2+s$. Conditional expectation and
\eqref{eq:third-Lambda} therefore give $\E Y(B)^3<\infty$.
Every bounded Borel set is contained in a ball, proving the third-moment
claim. In particular, the full second and third measures are locally
finite, so their densities in \eqref{eq:band-density-sum} are finite
almost everywhere globally.

\subsection{Finite fourth density away from a null exceptional set}
\label{app:radius-bound}

The fourth measure has infinite mass on some bounded sets, so a moment
bound cannot prove its density finite. Instead, we show that near almost
every quadruple, sufficiently large radius bands make no contribution.
Only finitely many locally finite fixed-band measures then remain.

Fix a noncoplanar quadruple $\mathbf x=(x_1,x_2,x_3,x_4)$, and translate
$x_1$ to zero. Put $u_i=x_i-x_1$, $i=2,3,4$, and let $W$ be the
invertible matrix with rows $u_i^{\mathsf T}$. Let $D_*$ denote the
quadruple diameter. For all sufficiently large $k$,
$2b_k+2M>D_*$. By \eqref{eq:surface-gap}, the four points then cannot
lie on spheres with two different centres. Any allocation in such a band
must use concentric spheres.

Let their common centre, relative to $x_1$, be $z$, and let $R=|z|$ be
the radius of the sphere through $x_1$. Because all candidate radii use
the same phase and step $H$, the other three radii have the form
$R+Hm_i$, where $m=(m_2,m_3,m_4)\in\mathbb Z^3$. The reverse triangle
inequality gives
\begin{equation}
 |Hm_i|\leq|u_i|\leq D_*.
 \label{eq:integer-bound}
\end{equation}
The equations for the other points are
$|u_i-z|^2=(R+Hm_i)^2$. Subtracting $|z|^2=R^2$ and putting $n=z/R$
gives
\begin{equation}
 Wn=-Hm+\frac b{2R},\qquad |n|=1,
 \qquad b_i=|u_i|^2-H^2m_i^2.
 \label{eq:concentric-equations}
\end{equation}
Taking norms after multiplication by $W^{-1}$ yields
\begin{equation}
 R\left|1-H|W^{-1}m|\right|
 \leq\frac12|W^{-1}b|.
 \label{eq:radius-bound}
\end{equation}
This bounds $R$ unless $H|W^{-1}m|=1$.

For $m=0$ the exceptional equality is impossible. For any fixed nonzero
integer vector $m$, it is contained in the zero set of the polynomial
\begin{equation}
 P_m(W)=H^2|\adj(W)m|^2-(\det W)^2.
 \label{eq:resonance-polynomial}
\end{equation}
This polynomial is not identically zero: substitution $W=tI$ gives
$t^4(H^2|m|^2-t^2)$. A nonzero real polynomial has a Lebesgue-null zero
set, as follows by induction on the number of variables and Fubini,
using the finitely many roots of a nonzero one-variable polynomial.
The coordinate change from the four vertices to $(x_1,W)$ is invertible
and linear. Consequently the union of all the sets $P_m(W)=0$, together
with the coplanarity set, is Lebesgue-null in the position space.
Call this exceptional union $\mathcal N$.

Take a quadruple outside $\mathcal N$. Choose a bounded neighbourhood on
which every diameter is at most $D_0$ and $W^{-1}$ remains bounded. Only
the finite integer set
\[
 \mathcal I=\{m\in\mathbb Z^3: |Hm_i|\leq D_0, i=2,3,4\}
\]
can appear in \eqref{eq:integer-bound}. At the chosen quadruple, the
numbers $|1-H|W^{-1}m||$, $m\in\mathcal I$, have a positive minimum.
By continuity and finiteness of $\mathcal I$, shrink the neighbourhood
so that all these numbers are at least a common $\varepsilon_0>0$.
The quantities $|W^{-1}b|$ have a common finite bound there, again because
$\mathcal I$ is finite. Equation~\eqref{eq:radius-bound} gives a uniform
upper bound on $R$ for every concentric allocation in that neighbourhood.
The other radii differ from $R$ by at most $D_0$, so they are bounded
as well.

For all sufficiently large bands, different-centre allocations are
excluded by \eqref{eq:surface-gap}, while concentric allocations are
excluded because the band's lower radius endpoint exceeds this common
bound. Thus, on the chosen bounded configuration neighbourhood, the full
fourth measure is a sum over only finitely many bands. Every such
fixed-band measure is locally finite by
Section~\ref{app:submersion}. Its sum has a density finite almost
everywhere on the neighbourhood.

The complement of $\mathcal N$ has a countable cover by such
neighbourhoods, so the fourth density is finite almost everywhere there.
The absolute-continuity proof in Section~\ref{app:submersion} already
shows that $\mathcal N$ itself has zero factorial measure. This finishes
the proof of Proposition~\ref{prop:sphere-regularity}, both inside and
outside the prescribed diameter domains.

\begin{thebibliography}{9}

\bibitem{Chapron2018}
A.~Chapron (2018).
\newblock A local Blaschke--Petkantschin formula in a Riemannian manifold.
\newblock arXiv:1807.07384v1.
\newblock \url{https://arxiv.org/abs/1807.07384}.

\bibitem{DVJ2008}
D.~J. Daley and D.~Vere-Jones (2008).
\newblock \emph{An Introduction to the Theory of Point Processes.
Volume II: General Theory and Structure}, second edition.
\newblock Springer, New York.
\newblock \url{https://doi.org/10.1007/978-0-387-49835-5}.

\bibitem{KaganKnopoff1980}
Y.~Y. Kagan and L.~Knopoff (1980).
\newblock Spatial distribution of earthquakes: the two-point correlation function.
\newblock \emph{Geophysical Journal of the Royal Astronomical Society}
\textbf{62}, 303--320.
\newblock \url{https://doi.org/10.1111/j.1365-246X.1980.tb04857.x}.

\bibitem{Kagan1981Three}
Y.~Y. Kagan (1981).
\newblock Spatial distribution of earthquakes: the three-point moment function.
\newblock \emph{Geophysical Journal of the Royal Astronomical Society}
\textbf{67}, 697--717.
\newblock \url{https://doi.org/10.1111/j.1365-246X.1981.tb06948.x}.

\bibitem{Kagan1981Four}
Y.~Y. Kagan (1981).
\newblock Spatial distribution of earthquakes: the four-point moment function.
\newblock \emph{Geophysical Journal of the Royal Astronomical Society}
\textbf{67}, 719--733.
\newblock \url{https://doi.org/10.1111/j.1365-246X.1981.tb06949.x}.

\end{thebibliography}
\end{document}
